\uuid{jXLf}
\titre{Système linéaire à paramètres/Méthode du Pivot de Gauss}
\auteur{hanani ; mzari}
\organisation{exo7}
\licence{CC BY-NC-SA 4.0}
\datecreate{2020-01-20}
\niveau{L1}
\module{Algèbre}
\chapitre{Déterminant, système linéaire}
\sousChapitre{Système linéaire, rang}
\difficulte{3}
\theme{Systèmes d'équations linéaires, Système linéaire à paramètres, Méthode du Pivot de Gauss}
\exo7qcmid{lille-2.systemes.010}

\contenu{
\begin{question}
\qclassification{200.03}
\qtags{motcle=Système linéaire à paramètres/Méthode du Pivot de Gauss}
On considère le système d'équations, d'inconnue $(x,y,z)\in \Rr^3$ et de paramètres des réels non nuls et distincts $a,b$ et $c$ :
$$(\mathtt{S}) \; \left\{\begin{array}{rcc}
ax+ay+bz&=&b\\
bx+by+cz&=&c\\
cx+cy+az&=&a.\end{array}\right.$$
Quelles sont les assertions vraies ?
\begin{answers}
\good{$(\mathtt{S}) \Leftrightarrow
\left\{\begin{array}{rcc}
ax+ay+bz&=&b\\
(ac-b^2)z&=&ac-b^2\\
(a^2-bc)z&=&a^2-bc.\end{array}\right.$}
\bad{$(\mathtt{S})$ n'admet pas de solution.}
\bad{$(\mathtt{S})$ admet une solution si et seulement si $a^2\neq bc$.}
\good{$(\mathtt{S})$ admet une infinité de solutions.}
\end{answers}
\begin{explanations} Comme $a,b$ et $c$ sont des réels non nuls,
$$(\mathtt{S}) \Leftrightarrow
\left\{\begin{array}{rcc}
ax+ay+bz&=&b\\
(ac-b^2)z&=&ac-b^2\\
(a^2-bc)z&=&a^2-bc.\\
\end{array}\right.$$
D'autre part, $a,b,c$ sont des réels distincts, on vérifie que $a^2\neq bc$ ou $b^2 \neq ac$. Par conséquent,
$(\mathtt{S})$ admet une infinité de solutions :
$\{(x,-x,1); \; x \in \Rr\}$.
\end{explanations}
\end{question}
}
